\begin{proof}[Proof of \cref{obj:thm:direct-reduction}]
Fix \(0<\beta\le1\). For integers \(n\ge2\), define
\[
r_{\mathrm{dir},n}:=n^{-\beta/(2\beta+1)}.
\]
Expectations of nonnegative losses and their suprema below are taken in the extended nonnegative reals.

\begin{enumerate}
\item \textbf{The selected radius at zero noise.}
By \cref{obj:lem:positive-endpoint}, there is a constant \(K_{\mathrm{end}}>0\), depending only on \(\beta\), such that, for every integer \(L\ge1\),
\[
M_{L,\beta}\le K_{\mathrm{end}}L^{-2\beta},
\qquad
\int_0^1 K_L(x)^2\,dx\le K_{\mathrm{end}}L^2.
\]
Define
\[
K_{\mathrm{cert}}
:=48K_{\mathrm{end}}+28\sqrt{30K_{\mathrm{end}}}>0.
\]
For a fixed \(n\ge2\), put
\[
a_{\mathrm{deg}}:=n^{1/(4\beta+2)},
\qquad
L_{\mathrm{dir}}:=\lfloor a_{\mathrm{deg}}\rfloor,
\qquad
L_{\max}:=\lceil a_{\mathrm{deg}}\rceil.
\]
Since \(a_{\mathrm{deg}}\ge1\), the floor inequalities give
\[
1\le L_{\mathrm{dir}},
\qquad
a_{\mathrm{deg}}<L_{\mathrm{dir}}+1\le2L_{\mathrm{dir}},
\qquad
L_{\mathrm{dir}}\le a_{\mathrm{deg}},
\qquad
L_{\mathrm{dir}}\le L_{\max}.
\]
These inequalities apply throughout \(a_{\mathrm{deg}}\ge1\), including \(1\le a_{\mathrm{deg}}<2\). Consequently,
\[
\begin{aligned}
M_{L_{\mathrm{dir}},\beta}
&\le K_{\mathrm{end}}L_{\mathrm{dir}}^{-2\beta}\\
&\le K_{\mathrm{end}}(a_{\mathrm{deg}}/2)^{-2\beta}\\
&=K_{\mathrm{end}}a_{\mathrm{deg}}^{-2\beta}2^{2\beta}
\le4K_{\mathrm{end}}a_{\mathrm{deg}}^{-2\beta},
\end{aligned}
\]
where the last inequality uses \(\beta\le1\).

In the variance formula of \cref{obj:synth_6}, the coefficient of the zeroth derivative is one at \(\sigma=0\), and every positive-order coefficient vanishes. Thus
\[
V_L(0)=\frac34\int_0^1K_L(x)^2\,dx,
\]
and hence
\[
\frac{40V_{L_{\mathrm{dir}}}(0)}{n}
\le\frac{30K_{\mathrm{end}}L_{\mathrm{dir}}^2}{n}
\le\frac{30K_{\mathrm{end}}a_{\mathrm{deg}}^2}{n}.
\]
Taking square roots yields
\[
\sqrt{\frac{40V_{L_{\mathrm{dir}}}(0)}{n}}
\le\sqrt{30K_{\mathrm{end}}}\,
\frac{a_{\mathrm{deg}}}{\sqrt n}.
\]
The two balancing identities are
\[
a_{\mathrm{deg}}^{-2\beta}
=n^{-2\beta/(4\beta+2)}
=r_{\mathrm{dir},n},
\qquad
\frac{a_{\mathrm{deg}}}{\sqrt n}
=n^{1/(4\beta+2)-1/2}
=r_{\mathrm{dir},n}.
\]
The positive-index radius in \cref{obj:synth_6} therefore satisfies
\[
\rho_{L_{\mathrm{dir}}}
=12M_{L_{\mathrm{dir}},\beta}
+28\sqrt{\frac{40V_{L_{\mathrm{dir}}}(0)}{n}}
\le K_{\mathrm{cert}}r_{\mathrm{dir},n}.
\]
Since \(L_{\mathrm{dir}}\) belongs to the candidate set, the minimization in \cref{obj:def:public-selection} gives
\[
E_\beta(n,0)=\rho_{L_\star}
\le\rho_{L_{\mathrm{dir}}}
\le K_{\mathrm{cert}}r_{\mathrm{dir},n}.
\]
% lean: CausalSmith.Stat.RdTruesideNoiseFrontier.certificate_zero_rate_bound

\item \textbf{Direct witnesses and their resolution.}
Fix \(n\ge2\) and \(\sigma\in[0,1]\), and define
\[
h_0:=n^{-1/(2\beta+1)}.
\]
Because \(n\ge1\) and \(2\beta+1>0\),
\[
0<h_0\le1,
\qquad
nh_0^{2\beta+1}=1,
\qquad
h_0^\beta=r_{\mathrm{dir},n}.
\]

For \(0<h\le1\), use the laws \(P_h^{\mathrm{dir},\pm}\) of \cref{obj:def:direct-alternatives}. Define their mean perturbation by
\[
\delta_h(x):=\kappa h^\beta v_h^{\mathrm{dir}}(x),
\qquad x\in[-1,1],
\qquad
\kappa=\frac1{100}.
\]
The profile bounds in \cref{obj:synth_7} imply
\[
0\le\delta_h(x)\le\kappa,
\qquad
|\delta_h(x)-\delta_h(t)|
\le\kappa|x-t|^\beta
\quad(x,t\in[-1,1]).
\]
Together with the conditional means in \cref{obj:def:direct-alternatives}, these give
\[
\frac14
\le\mu_d(P_h^{\mathrm{dir},\pm},x)
\le\frac34,
\qquad
[\mu_0(P_h^{\mathrm{dir},\pm},\cdot)]_\beta=0,
\qquad
[\mu_1(P_h^{\mathrm{dir},\pm},\cdot)]_\beta\le\kappa\le2.
\]
Their score density is continuous and satisfies
\[
f(P_h^{\mathrm{dir},\pm},x)=\frac12
\quad(x\in[-1,1]),
\qquad
\int_{-1}^1f(P_h^{\mathrm{dir},\pm},x)\,dx=1.
\]
Their Gaussian error and independence properties are part of their construction. Thus all the conditions of \cref{obj:def:model} hold:
\[
P_h^{\mathrm{dir},\pm}\in\mathcal M_\beta(\sigma)
\quad(0<h\le1,\ 0\le\sigma\le1).
\]
Since \(v_h^{\mathrm{dir}}(0)=1\), their targets satisfy
\[
\theta(P_h^{\mathrm{dir},+})=\kappa h^\beta,
\qquad
\theta(P_h^{\mathrm{dir},-})=-\kappa h^\beta,
\qquad
\left|\theta(P_h^{\mathrm{dir},+})
-\theta(P_h^{\mathrm{dir},-})\right|
=2\kappa h^\beta.
\]
% lean: CausalSmith.Stat.RdTruesideNoiseFrontier.directAltLaw_model
% lean: CausalSmith.Stat.RdTruesideNoiseFrontier.directResolution_balance
% lean: CausalSmith.Stat.RdTruesideNoiseFrontier.directAltLaw_separation

\item \textbf{Noiseless information bound.}
For \(0<h\le1\), define the one-observation laws
\[
\mathsf P_{\sigma,h}^{\pm}
:=\operatorname{Law}_{P_h^{\mathrm{dir},\pm}}
(X+\sigma\varepsilon,D,Y).
\]
For probability measures \(\mathsf Q_1,\mathsf Q_2\) on a common measurable space, use
\[
\operatorname{TV}(\mathsf Q_1,\mathsf Q_2)
:=\sup_{\mathsf F\ \mathrm{measurable}}
|\mathsf Q_1(\mathsf F)-\mathsf Q_2(\mathsf F)|.
\]
In the absolutely continuous setting used here, the chi-squared divergence is
\[
\chi^2(\mathsf Q_1,\mathsf Q_2)
:=\int
\left(\frac{d\mathsf Q_1}{d\mathsf Q_2}-1\right)^2
\,d\mathsf Q_2.
\]

Let the reference measure on \(\mathbb R\times\{0,1\}^2\) be defined, for every measurable \(\mathsf F\), by
\[
\eta_{\mathrm{obs}}(\mathsf F)
:=\sum_{d=0}^1\sum_{y=0}^1
\int_{\mathbb R}\mathbf1_{\mathsf F}(x,d,y)\,dx.
\]
The densities
\[
p_h^\pm:=\frac{d\mathsf P_{0,h}^\pm}{d\eta_{\mathrm{obs}}}
\]
are given by
\[
\begin{aligned}
p_h^\pm(x,1,1)&=\frac14\pm\frac{\delta_h(x)}2,
&&0\le x\le1,\\
p_h^\pm(x,1,0)&=\frac14\mp\frac{\delta_h(x)}2,
&&0\le x\le1,\\
p_h^\pm(x,0,y)&=\frac14,
&&-1\le x<0,\quad y\in\{0,1\},
\end{aligned}
\]
and are zero elsewhere. Define a likelihood ratio on the entire observation space by
\[
\Lambda_h(x,d,y):=
\begin{cases}
p_h^+(x,1,y)/p_h^-(x,1,y),
&d=1,\ 0\le x\le1,\\
1,&\text{otherwise}.
\end{cases}
\]
The treated-cell densities lie between \(1/8\) and \(3/8\), so
\[
d\mathsf P_{0,h}^+
=\Lambda_h\,d\mathsf P_{0,h}^-,
\qquad
0\le\Lambda_h\le3,
\qquad
(\Lambda_h-1)^2\le4.
\]
In particular,
\[
\mathsf P_{0,h}^+\ll\mathsf P_{0,h}^-,
\qquad
\int(\Lambda_h-1)^2\,d\mathsf P_{0,h}^-<\infty.
\]

For each treated cell,
\[
p_h^+(x,1,y)-p_h^-(x,1,y)
=(2y-1)\delta_h(x),
\qquad
p_h^-(x,1,y)\ge\frac18.
\]
Consequently,
\[
0\le
p_h^-(x,1,y)(\Lambda_h(x,1,y)-1)^2
=\frac{\delta_h(x)^2}{p_h^-(x,1,y)}
\le8\delta_h(x)^2.
\]
Outside the treated cells, the corresponding integrand is zero. Summing over the two outcome cells gives
\[
\chi^2(\mathsf P_{0,h}^+,\mathsf P_{0,h}^-)
\le16\kappa^2h^{2\beta}
\int_0^1[v_h^{\mathrm{dir}}(x)]^2\,dx.
\]
The profile is \(1-x/h\) on \([0,h]\) and zero on \([h,1]\). With the substitution \(t=1-x/h\),
\[
\int_0^1[v_h^{\mathrm{dir}}(x)]^2\,dx
=\int_0^h(1-x/h)^2\,dx
=h\int_0^1t^2\,dt
=\frac h3.
\]
It follows that
\[
\chi^2(\mathsf P_{0,h}^+,\mathsf P_{0,h}^-)
\le\frac{16}{3}\kappa^2h^{2\beta+1}.
\]
At \(h=h_0\), the balance identity therefore yields
\[
n\chi^2(\mathsf P_{0,h_0}^+,\mathsf P_{0,h_0}^-)
\le\frac{16}{3}\kappa^2
=\frac{16}{30000}
\le\frac1{100}.
\]
% lean: CausalSmith.Stat.RdTruesideNoiseFrontier.directAltLaw_observed_ac_zero
% lean: CausalSmith.Stat.RdTruesideNoiseFrontier.directAltLaw_observed_likelihood_square_integrable_zero
% lean: CausalSmith.Stat.RdTruesideNoiseFrontier.directAltLaw_observed_chiSqDiv_zero_le

\item \textbf{Product testing bound and Gaussian contraction.}
Define the sample laws by
\[
\mathsf S_{\sigma,h}^{\pm}
:=(\mathsf P_{\sigma,h}^{\pm})^{\otimes n},
\]
and abbreviate the one-observation divergence at the selected resolution as
\[
\chi_{\mathrm{dir}}
:=\chi^2(\mathsf P_{0,h_0}^+,\mathsf P_{0,h_0}^-).
\]
The product likelihood ratio is
\[
\frac{d\mathsf S_{0,h_0}^+}{d\mathsf S_{0,h_0}^-}
\bigl((o_i)_{i=1}^n\bigr)
=\prod_{i=1}^n\Lambda_{h_0}(o_i),
\qquad
(o_i)_{i=1}^n\in(\mathbb R\times\{0,1\}^2)^n,
\]
almost everywhere under \(\mathsf S_{0,h_0}^-\). Since the one-observation likelihood ratio integrates to one,
\[
\int\Lambda_{h_0}\,d\mathsf P_{0,h_0}^-=1,
\qquad
\chi_{\mathrm{dir}}
=\int(\Lambda_{h_0}-1)^2\,d\mathsf P_{0,h_0}^-
=\int\Lambda_{h_0}^2\,d\mathsf P_{0,h_0}^- -1.
\]
Product integration likewise gives
\[
\int\prod_{i=1}^n\Lambda_{h_0}(o_i)
\,d\mathsf S_{0,h_0}^-\bigl((o_i)_{i=1}^n\bigr)
=\left(\int\Lambda_{h_0}\,d\mathsf P_{0,h_0}^-\right)^n
=1.
\]
Expanding the squared likelihood deviation and factoring its second moment over the product measure therefore yields
\[
\begin{aligned}
1+\chi^2(\mathsf S_{0,h_0}^+,\mathsf S_{0,h_0}^-)
&=\int\left(\prod_{i=1}^n\Lambda_{h_0}(o_i)\right)^2
\,d\mathsf S_{0,h_0}^-\bigl((o_i)_{i=1}^n\bigr)\\
&=\left(\int\Lambda_{h_0}^2\,d\mathsf P_{0,h_0}^-\right)^n\\
&=(1+\chi_{\mathrm{dir}})^n.
\end{aligned}
\]
Since \(\chi_{\mathrm{dir}}\ge0\) and \(n\chi_{\mathrm{dir}}\le1/100\),
\[
(1+\chi_{\mathrm{dir}})^n
\le e^{n\chi_{\mathrm{dir}}}
\le e^{1/100}
\le\sum_{j=0}^\infty(1/100)^j
=\frac{100}{99}.
\]
Here the penultimate inequality follows from the exponential series and \(j!\ge1\). Thus
\[
\chi^2(\mathsf S_{0,h_0}^+,\mathsf S_{0,h_0}^-)
\le\frac1{99}\le\frac4{25}.
\]
For absolutely continuous probability measures, the signed density difference has integral zero, and Cauchy--Schwarz gives
\[
\operatorname{TV}(\mathsf Q_1,\mathsf Q_2)
=\frac12\int
\left|\frac{d\mathsf Q_1}{d\mathsf Q_2}-1\right|
\,d\mathsf Q_2
\le\frac12\sqrt{\chi^2(\mathsf Q_1,\mathsf Q_2)}.
\]
Applying this to the product laws gives
\[
\operatorname{TV}(\mathsf S_{0,h_0}^+,\mathsf S_{0,h_0}^-)
\le\frac12\sqrt{\frac4{25}}=\frac15.
\]

For a sample
\[
\mathbf s=((x_i,d_i,y_i))_{i=1}^n
\in(\mathbb R\times\{0,1\}^2)^n,
\]
define the common Gaussian observation kernel by
\[
\mathsf K_{\sigma,n}(\mathbf s,\mathsf F)
:=
\Pr\!\left\{
((x_i+\sigma\varepsilon_i,d_i,y_i))_{i=1}^n
\in\mathsf F
\right\},
\qquad
(\varepsilon_i)_{i=1}^n\sim N(0,1)^{\otimes n}.
\]
The Gaussian variables in this definition are independent of the input sample. Gaussian error independence under the direct laws implies
\[
\mathsf S_{\sigma,h_0}^{\pm}
=\mathsf K_{\sigma,n}\mathsf S_{0,h_0}^{\pm}.
\]
For every measurable output event \(\mathsf F\), the function
\(\mathsf K_{\sigma,n}(\cdot,\mathsf F)\) takes values in \([0,1]\). The layer-cake formula therefore bounds the difference of its integrals by
\[
\begin{aligned}
&\left|
\int\mathsf K_{\sigma,n}(\mathbf s,\mathsf F)
\,d\mathsf S_{0,h_0}^+(\mathbf s)
-
\int\mathsf K_{\sigma,n}(\mathbf s,\mathsf F)
\,d\mathsf S_{0,h_0}^-(\mathbf s)
\right|\\
&\qquad\le
\int_0^1
\left|
\mathsf S_{0,h_0}^+
\{\mathsf K_{\sigma,n}(\cdot,\mathsf F)>t\}
-
\mathsf S_{0,h_0}^-
\{\mathsf K_{\sigma,n}(\cdot,\mathsf F)>t\}
\right|\,dt\\
&\qquad\le
\operatorname{TV}(\mathsf S_{0,h_0}^+,\mathsf S_{0,h_0}^-).
\end{aligned}
\]
Taking the supremum over output events proves
\[
\operatorname{TV}(\mathsf S_{\sigma,h_0}^+,\mathsf S_{\sigma,h_0}^-)
\le
\operatorname{TV}(\mathsf S_{0,h_0}^+,\mathsf S_{0,h_0}^-)
\le\frac15.
\]
The construction and this inequality also apply at \(\sigma=0\).
% lean: Causalean.Stat.iid_tv_le_one_fifth_of_chiSqDiv
% lean: CausalSmith.Stat.RdTruesideNoiseFrontier.sampleLaw_tv_le_zero_of_model

\item \textbf{Transfer to estimation risk and honest length.}
Let
\[
\mathsf U_{\mathrm{seed}}:=\operatorname{Uniform}[0,1],
\qquad
\mathsf E_{\sigma,h_0}^{\pm}
:=\mathsf S_{\sigma,h_0}^{\pm}\otimes\mathsf U_{\mathrm{seed}}.
\]
The product total-variation inequality gives
\[
\begin{aligned}
\operatorname{TV}(\mathsf E_{\sigma,h_0}^+,\mathsf E_{\sigma,h_0}^-)
&\le
\operatorname{TV}(\mathsf S_{\sigma,h_0}^+,\mathsf S_{\sigma,h_0}^-)
+\operatorname{TV}(\mathsf U_{\mathrm{seed}},\mathsf U_{\mathrm{seed}})\\
&\le\frac15.
\end{aligned}
\]
Define the common submeasure by
\[
d\mathsf H_{\sigma,h_0}
:=
\min\!\left\{
\frac{d\mathsf E_{\sigma,h_0}^+}
{d(\mathsf E_{\sigma,h_0}^++\mathsf E_{\sigma,h_0}^-)},
\frac{d\mathsf E_{\sigma,h_0}^-}
{d(\mathsf E_{\sigma,h_0}^++\mathsf E_{\sigma,h_0}^-)}
\right\}
d(\mathsf E_{\sigma,h_0}^++\mathsf E_{\sigma,h_0}^-).
\]
It satisfies
\[
\mathsf H_{\sigma,h_0}\le\mathsf E_{\sigma,h_0}^{\pm},
\qquad
\int d\mathsf H_{\sigma,h_0}
=
1-\operatorname{TV}(\mathsf E_{\sigma,h_0}^+,\mathsf E_{\sigma,h_0}^-)
\ge\frac45.
\]
The mass identity follows from
\(\min\{u,v\}=(u+v-|u-v|)/2\).

For any admissible estimator \(\widehat\theta\), domination by the common submeasure and the triangle inequality imply
\[
\begin{aligned}
&\mathbb E_{P_{h_0}^{\mathrm{dir},+},n}
|\widehat\theta-\theta(P_{h_0}^{\mathrm{dir},+})|
+
\mathbb E_{P_{h_0}^{\mathrm{dir},-},n}
|\widehat\theta-\theta(P_{h_0}^{\mathrm{dir},-})|\\
&\quad\ge
\int\left(
|\widehat\theta-\theta(P_{h_0}^{\mathrm{dir},+})|
+
|\widehat\theta-\theta(P_{h_0}^{\mathrm{dir},-})|
\right)\,d\mathsf H_{\sigma,h_0}\\
&\quad\ge
2\kappa h_0^\beta\int d\mathsf H_{\sigma,h_0}
\ge\frac85\kappa h_0^\beta.
\end{aligned}
\]
At least one of the two extended risks is therefore at least
\(\frac45\kappa h_0^\beta\), and both laws belong to the model class. Hence
\[
\inf_{\widehat\theta}
\sup_{P\in\mathcal M_\beta(\sigma)}
\mathbb E_{P,n}|\widehat\theta-\theta(P)|
\ge\frac45\kappa h_0^\beta.
\]

For any \(I\in\mathcal I_{\beta,n,\sigma}\), define its two coverage events by
\[
\mathsf F_{h_0,I}^{\pm}
:=
\{(S_n,U):\theta(P_{h_0}^{\mathrm{dir},\pm})\in I(S_n,U)\}.
\]
Honesty as defined in \cref{obj:def:honest-class} gives
\[
\mathsf E_{\sigma,h_0}^{+}(\mathsf F_{h_0,I}^{+})\ge\frac9{10},
\qquad
\mathsf E_{\sigma,h_0}^{-}(\mathsf F_{h_0,I}^{-})\ge\frac9{10}.
\]
Total variation transfers the first probability to the minus law:
\[
\mathsf E_{\sigma,h_0}^{-}(\mathsf F_{h_0,I}^{+})
\ge\frac9{10}-\frac15=\frac7{10}.
\]
Using the union bound under that law,
\[
\mathsf E_{\sigma,h_0}^{-}
(\mathsf F_{h_0,I}^{+}\cap\mathsf F_{h_0,I}^{-})
\ge\frac7{10}+\frac9{10}-1=\frac35.
\]
On the intersection, the connected interval contains both targets and has length at least their separation. Thus
\[
\mathbb E_{P_{h_0}^{\mathrm{dir},-},n}\operatorname{length}(I)
\ge2\kappa h_0^\beta
\mathsf E_{\sigma,h_0}^{-}
(\mathsf F_{h_0,I}^{+}\cap\mathsf F_{h_0,I}^{-})
\ge\frac65\kappa h_0^\beta.
\]
Taking the infimum over honest intervals yields
\[
\inf_{I\in\mathcal I_{\beta,n,\sigma}}
\sup_{P\in\mathcal M_\beta(\sigma)}
\mathbb E_{P,n}\operatorname{length}(I)
\ge\frac65\kappa h_0^\beta.
\]
% lean: CausalSmith.Stat.RdTruesideNoiseFrontier.experiment_tv_le
% lean: CausalSmith.Stat.RdTruesideNoiseFrontier.twoPoint_absRisk_sum
% lean: CausalSmith.Stat.RdTruesideNoiseFrontier.twoPoint_interval_lower
% lean: CausalSmith.Stat.RdTruesideNoiseFrontier.direct_two_point_transfer

\item \textbf{Real lower bounds.}
For every \(P\in\mathcal M_\beta(\sigma)\), the mean bounds give
\[
|\theta(P)|
=|\mu_1(P,0)-\mu_0(P,0)|
\le\frac12.
\]
The constant-zero estimator consequently satisfies
\[
\sup_{P\in\mathcal M_\beta(\sigma)}
\mathbb E_{P,n}|0-\theta(P)|\le\frac12.
\]
The constant interval \([-1,1]\) satisfies
\[
\mathbb P_{P,n}\{\theta(P)\in[-1,1]\}=1,
\qquad
\mathbb E_{P,n}\operatorname{length}([-1,1])=2
\quad(P\in\mathcal M_\beta(\sigma)).
\]
It is therefore an admissible honest interval. These two procedures bound the extended minimax values by finite numbers, so conversion to real numbers preserves the preceding lower inequalities. Since \(h_0^\beta=r_{\mathrm{dir},n}\),
\[
R_\beta(n,\sigma)\ge\frac45\kappa r_{\mathrm{dir},n},
\qquad
\mathcal L_\beta(n,\sigma)\ge\frac65\kappa r_{\mathrm{dir},n}.
\]
Define
\[
c:=\frac45\kappa>0.
\]
As \(\kappa r_{\mathrm{dir},n}\ge0\), both criteria are at least
\(c\,r_{\mathrm{dir},n}\), uniformly over \(n\ge2\) and \(\sigma\in[0,1]\).
% lean: CausalSmith.Stat.RdTruesideNoiseFrontier.risk_value_ne_top
% lean: CausalSmith.Stat.RdTruesideNoiseFrontier.length_value_ne_top
% lean: CausalSmith.Stat.RdTruesideNoiseFrontier.direct_uniform_lower

\item \textbf{Attainment and minimax upper bounds.}
Set
\[
C:=2K_{\mathrm{cert}}>0,
\qquad
P_{\mathrm{unit}}:=P_1^{\mathrm{dir},+}.
\]
The witness verification above, at \(h=1\), gives
\[
P_{\mathrm{unit}}\in\mathcal M_\beta(0).
\]
For every \(n\ge2\), apply \cref{obj:thm:finite-certificate} with
\(\sigma=0\), \(L=1\), and \(P=P_{\mathrm{unit}}\). Its selected-procedure conclusions give
\[
I_{L_\star}\in\mathcal I_{\beta,n,0},
\]
\[
\sup_{P\in\mathcal M_\beta(0)}
\mathbb E_{P,n}
|\widehat\theta_{L_\star}-\theta(P)|
\le E_\beta(n,0),
\qquad
\sup_{P\in\mathcal M_\beta(0)}
\mathbb E_{P,n}\operatorname{length}(I_{L_\star})
\le2E_\beta(n,0).
\]
Since \(K_{\mathrm{cert}}r_{\mathrm{dir},n}\ge0\), the radius bound from the first step implies
\[
\sup_{P\in\mathcal M_\beta(0)}
\mathbb E_{P,n}
|\widehat\theta_{L_\star}-\theta(P)|
\le K_{\mathrm{cert}}r_{\mathrm{dir},n}
\le C r_{\mathrm{dir},n},
\]
\[
\sup_{P\in\mathcal M_\beta(0)}
\mathbb E_{P,n}\operatorname{length}(I_{L_\star})
\le2K_{\mathrm{cert}}r_{\mathrm{dir},n}
=C r_{\mathrm{dir},n}.
\]
Each extended minimax value is bounded above by the worst-case loss of its displayed admissible procedure. The selected interval's honesty supplies its admissibility for the length minimization. The finite upper bounds permit order-preserving conversion to real numbers, giving
\[
R_\beta(n,0)\le C r_{\mathrm{dir},n},
\qquad
\mathcal L_\beta(n,0)\le C r_{\mathrm{dir},n}.
\]
The constants \(c\) and \(C\) depend only on \(\beta\), and the displayed inequalities and honesty statement hold for every integer \(n\ge2\), as required.
% lean: CausalSmith.Stat.RdTruesideNoiseFrontier.risk_le_of_attainerRisk_le
% lean: CausalSmith.Stat.RdTruesideNoiseFrontier.lengthRisk_le_of_attainerLength_le
% lean: CausalSmith.Stat.RdTruesideNoiseFrontier.direct_reduction
\end{enumerate}
\end{proof}
